% ===== PRÉAMBULE CANONIQUE — à coller VERBATIM en tête de CHAQUE .tex =====
\documentclass[11pt,a4paper]{article}
\usepackage[utf8]{inputenc}\usepackage[T1]{fontenc}\usepackage{lmodern}
\usepackage[french]{babel}
\usepackage{amsmath,amssymb,mathtools}\usepackage{icomma}
\usepackage[version=4]{mhchem}          % \ce{...} : équations & formules
\usepackage{siunitx}\sisetup{locale=FR} % \qty{}{}, \si{}, \ang{}
\usepackage{chemfig}                    % \chemfig{...} : structures développées
\usepackage{xcolor}\usepackage{colortbl}
\usepackage{tikz}\usepackage{pgfplots}\pgfplotsset{compat=1.18}
\usepackage[most]{tcolorbox}\usepackage{enumitem}\usepackage{array,booktabs}
\usepackage[a4paper,margin=2.2cm]{geometry}
\usetikzlibrary{arrows.meta,calc,angles,quotes,positioning,patterns,backgrounds,shapes.geometric}
\definecolor{primary}{RGB}{222,49,99}\definecolor{primarydark}{RGB}{150,22,55}
\definecolor{defcol}{RGB}{0,114,178}\definecolor{methcol}{RGB}{52,92,148}
\definecolor{excol}{RGB}{0,135,90}\definecolor{attncol}{RGB}{200,40,55}
\tcbset{boxbase/.style={enhanced,breakable,boxrule=0.9pt,arc=2.5pt,
  left=3mm,right=3mm,top=2.3mm,bottom=2.3mm,fonttitle=\bfseries,coltitle=white}}
% --- boîtes TUTORIEL (noms EXACTS, ne pas renommer) ---
\newtcolorbox{cle}[1][]{boxbase,colback=defcol!6,colframe=defcol,
  title={\textsc{À retenir}\ifx\relax#1\relax\relax\else\textmd{ : \itshape #1}\fi}}
\newtcolorbox{methode}[1][]{boxbase,colback=methcol!7,colframe=methcol,sharp corners,
  title={\textsc{Méthode}\ifx\relax#1\relax\relax\else\textmd{ --- \itshape #1}\fi}}
\newtcolorbox{exemplebox}[1][]{boxbase,colback=excol!5,colframe=excol,
  title={\textsc{Exemple}\ifx\relax#1\relax\relax\else\textmd{ : \itshape #1}\fi}}
\newtcolorbox{piege}[1][]{boxbase,colback=attncol!6,colframe=attncol,
  title={\textsc{Piège}\ifx\relax#1\relax\relax\else\textmd{ : \itshape #1}\fi}}
\newtcolorbox{remarque}[1][]{boxbase,colback=black!4,colframe=black!45,
  title={\textsc{Remarque}\ifx\relax#1\relax\relax\else\textmd{ : \itshape #1}\fi}}
% --- boîtes EXERCICE / CORRIGÉ (noms EXACTS, auto-comptées) ---
\newtcolorbox[auto counter]{exo}[1][]{boxbase,colback=primary!4,colframe=primary,
  title={\textsc{Exercice}~\thetcbcounter\ifx\relax#1\relax\relax\else\textmd{ --- \itshape #1}\fi}}
\newtcolorbox[auto counter]{corrige}[1][]{boxbase,colback=excol!4,colframe=excol,
  title={\textsc{Corrigé}~\thetcbcounter\ifx\relax#1\relax\relax\else\textmd{ --- \itshape #1}\fi}}
\newcommand{\R}{\mathbb{R}}\newcommand{\N}{\mathbb{N}}\newcommand{\Z}{\mathbb{Z}}\newcommand{\ds}{\displaystyle}
% (un \usepackage{fancyhdr}+en-tête est possible ; il est ignoré côté web)
% =========================================================================
\begin{document}

\begin{exo}[le double piège du cuivre et de l'azote]
Les valences variables et les préfixes fractionnaires sont sources d'erreurs classiques.
\begin{enumerate}[label=\alph*)]
  \item Nomme \ce{CuO} et \ce{Cu2O}.
  \item Nomme \ce{NO} et \ce{NO2}.
  \item Un étudiant nomme \ce{N2O} « dioxyde d'azote ». Est-ce correct ? Sinon, donne le nom de \ce{N2O}
        (convention du rapport O/X) et la formule réellement nommée « dioxyde d'azote ».
\end{enumerate}
\end{exo}

\begin{exo}[repérer le nom incorrect]
Pour chaque proposition, dis si le nom est correct ; sinon corrige-le.
\begin{enumerate}[label=\alph*)]
  \item \ce{Cu2O} nommé « oxyde de cuivre(II) »
  \item \ce{Fe2O3} nommé « oxyde de fer(III) »
  \item \ce{N2O5} nommé « pentoxyde d'azote »
\end{enumerate}
\end{exo}

\begin{exo}[de l'hydroxyde à l'oxyde métallique]
À un hydroxyde correspond un oxyde du \emph{même} métal, à la \emph{même} valence (par déshydratation :
$\ce{2 M(OH)_n -> M2O_n + $n$ H2O}$). Donne la formule et le nom de l'oxyde associé.
\begin{enumerate}[label=\alph*)]
  \item hydroxyde de fer(III) \ce{Fe(OH)3}
  \item hydroxyde de calcium \ce{Ca(OH)2}
  \item hydroxyde d'aluminium \ce{Al(OH)3}
\end{enumerate}
\end{exo}

\begin{exo}[l'oxyde acide et son oxacide (anhydride)]
Un oxyde non métallique est l'\emph{anhydride} d'un oxacide : le non-métal y garde le \emph{même} n.o.
\textbf{Données :} acide sulfurique \ce{H2SO4} (soufre au degré $+\mathrm{VI}$), acide sulfureux \ce{H2SO3}
(soufre au degré $+\mathrm{IV}$).
\begin{enumerate}[label=\alph*)]
  \item Détermine le n.o.\ du soufre dans \ce{SO3}, puis nomme l'oxacide correspondant.
  \item Même travail pour \ce{SO2}.
  \item La réaction $\ce{SO3 + H2O -> H2SO4}$ change-t-elle le n.o.\ du soufre ? Justifie.
\end{enumerate}
\end{exo}

\begin{exo}[synthèse : nom $\to$ formule $\to$ n.o.]
On considère l'\textbf{hémipentoxyde de phosphore}.
\begin{enumerate}[label=\alph*)]
  \item Écris sa formule brute.
  \item Détermine le n.o.\ du phosphore dans ce composé.
  \item Vérifie la neutralité électrique de la molécule.
\end{enumerate}
\end{exo}

\end{document}
